\section{Detalle del Conector}

\begin{figure}[h!]
    \centering
    % TODO
    % \includegraphics[width=0.8\textwidth]{path_to_the_image/imagedetail.png}
    \caption{Detalle ``A'' del Conector REEL-Contenedor y Acople del Chicote}
    \label{fig:detalle_a}
\end{figure}

El componente mostrado en la Figura \ref{fig:detalle_a} consiste en la parte de acoplamiento perteneciente al sistema de conectores modelo REEL. El detalle ``A'' especifica la estructura interna de las conexiones entre el ``Chicote'' y el cuerpo principal del conector, identificado aquí como `COPL'.

\begin{itemize}
    \item \textbf{Chicote:} Segmento flexible utilizado para la transmisión de movimientos mecánicos y/o eléctricos a través de su estructura articulada.
    \item \textbf{Conector REEL-Contenedor:} Este elemento se encarga de acoger las terminaciones del chicote y asegurar la transmisión eficaz de los componentes transportados, ya sea mecánica o eléctricamente.
    \item \textbf{Acople (COPL):} Parte del conector que facilita la conexión entre el chicote y el sistema al que este será integrado. Permite una acoplamiento seguro y estable.
\end{itemize}

Es crucial asegurarse de que todas las uniones entre estos componentes estén bien selladas y aseguradas para evitar cualquier tipo de fuga o desconexión imprevista que podría comprometer la funcionalidad del sistema. Estos detalles son fundamentales para mantener la integridad y funcionalidad en entornos que requieren altos estándares de seguridad y eficiencia, como es el caso en aplicaciones de inspección y manejo radiactivo donde se utilizan este tipo de componentes.
